% ============================================================
%  Главна датотека — уноси поглавља из директоријума chapters/
%  Компајлирање: pdflatex seminar.tex  (покренути два пута)
% ============================================================
\documentclass[12pt, a4paper]{article}

% ---------- Пакети ----------
\usepackage[utf8]{inputenc}
\usepackage[T2A]{fontenc}
\usepackage[serbianc]{babel}
\usepackage{amsmath, amssymb, amsthm}
\usepackage{graphicx}
\usepackage{booktabs}
\usepackage{hyperref}
\usepackage{geometry}
\usepackage{fancyhdr}
\usepackage{titling}
\usepackage{setspace}
\usepackage{parskip}
\usepackage{xcolor}
\usepackage{caption}
\usepackage{enumitem}
\usepackage{tcolorbox}
\usepackage{algpseudocode}
\usepackage{algorithm}
\usepackage{tikz}
\usetikzlibrary{graphs, positioning}
\usepackage{qcircuit}   % квантна кола (опционо, уклони ако није инсталирано)

% ---------- Геометрија ----------
\geometry{a4paper, top=3cm, bottom=3cm, left=3cm, right=3cm}

% ---------- Проред ----------
\setstretch{1.15}
\setlength{\parindent}{0pt}
\setlength{\parskip}{8pt}

% ---------- Заглавље / подножје ----------
\pagestyle{fancy}
\fancyhf{}
\fancyhead[L]{\small\textit{Семинар катедре за рачунарство и информатику}}
\fancyhead[R]{\small\textit{КР и НП-тешки проблеми}}
\fancyfoot[C]{\small\thepage}
\renewcommand{\headrulewidth}{0.4pt}
\renewcommand{\footrulewidth}{0pt}

% ---------- Теоремска окружења ----------
\newtheorem{theorem}{Теорема}[section]
\newtheorem{lemma}[theorem]{Лема}
\newtheorem{definition}[theorem]{Дефиниција}
\newtheorem{example}[theorem]{Пример}
\theoremstyle{remark}
\newtheorem*{remark}{Напомена}

% ---------- Hyperref боје ----------
\hypersetup{colorlinks=true, linkcolor=black, citecolor=black, urlcolor=blue!60!black}

% ---------- Кутија са информацијама ----------
\tcbset{
  infostyle/.style={
    colback=gray!8, colframe=gray!40,
    boxrule=0.4pt, arc=2pt,
    left=6pt, right=6pt, top=4pt, bottom=4pt,
    fonttitle=\bfseries\small
  }
}

% ============================================================
%  БАНЕР КОНФЕРЕНЦИЈЕ
% ============================================================
\newcommand{\conferencebanner}{%
  \begin{tcolorbox}[infostyle, title=Информације]
    \small
    \textbf{Догађај:} Семинар катедре за рачунарство и информатику
    \textbf{Датум:} 12. април 2026. \quad
    \textbf{Место:} Математички факултет, Универзитет у Београду
    \textbf{Тема:} Квантно рачунарство и примене
  \end{tcolorbox}
  \vspace{0.6cm}
}

% ============================================================
%  НАСЛОВНА СТРАНА
% ============================================================
\title{%
  {\large\scshape Семинар катедре за рачунарство}\\[8pt]
  \rule{\linewidth}{1.2pt}\\[10pt]
  {\LARGE\bfseries Квантно рачунарство и НП-тешки проблеми:\\[4pt]
   \Large Од MaxCut проблема до аутоматског резоновања}\\[6pt]
  \rule{\linewidth}{0.5pt}
}

\author{%
  \textbf{Александар Ивановић}\thanks{Аутор за контакт.
    Контакт: \texttt{aleksandar.ivke96@gmail.com}}\\
  \small Катедра за рачунарство, Математички факултет\\
  \small Универзитет у Београду\\[6pt]
}

\date{12. април 2026.}

% ============================================================
\begin{document}
% ============================================================

\maketitle
\thispagestyle{fancy}
\vspace{0.4cm}
\conferencebanner

% ============================================================
%  САЖЕТАК
% ============================================================
\begin{abstract}
\noindent
Овај рад истражује потенцијал квантног рачунарства у решавању НП-тешких
проблема, са акцентом на проблем MaxCut као репрезентативном примеру.
Кроз призму класичног аутоматског резоновања и Z3 решавача, уводимо читаоца
у ограничења класичних приступа, а затим показујемо како квантни алгоритми
— посебно QAOA (Quantum Approximate Optimization Algorithm) — нуде нову
парадигму за решавање ових проблема.
Рад повезује формалне методе верификације са квантним рачунарством и указује
на правце будућег истраживања.
\end{abstract}

\textbf{Кључне речи:} квантно рачунарство, НП-тешки проблеми, MaxCut,
Z3 решавач, QAOA, аутоматско резоновање.

\vspace{0.5em}
\hrule
\vspace{1.2em}

\tableofcontents
\newpage

% ============================================================
%  ПОГЛАВЉА
% ============================================================
% ============================================================
%  Поглавље 1: Увод
% ============================================================
\section{Увод}
\label{sec:uvod}

Рачунарство је суочено са класом проблема које класични рачунари
не могу да решавају ефикасно за велике улазе, такозваним НП-тешким
проблемима. Ови проблеми јављају се у логистици, криптографији,
машинском учењу и верификацији софтвера, па је њихово ефикасно решавање
од велике практичне важности.


Квантно рачунарство нуди нову рачунску парадигму засновану на принципима
квантне механике, суперпозицији и спрегнутости, која потенцијално
омогућује истовремено претраживање експоненцијално великог простора стања.
Квантна надмоћ над класичним рачунарима за НП-тешке проблеме
теоријски није доказана, али хибридни квантно-класични алгоритми показују
значајан практичан потенцијал.

У овом раду ћемо користи проблем \textit{MaxCut} као водећи пример кроз који
ћемо покушати да приближимо могуће предности једне гране квантног рачунарства.
Представићемо модел \textbf{квантног жарења (Quantum Annealing)} и одговарајући математички формализам 
 \textbf{квадратне неограничене бинарне оптимизације (QUBO quadratic unconstrained binary optimisation)}
и на њему моделовати проблем \textit{MaxCut}. Такође ћемо исти проблем моделовати
коришћењем \textit{Z3} решавача и указати на везу са \textit{QUBO}.
% ============================================================
%  Поглавље 2: Квантно рачунарство и две главне парадигме
% ============================================================
\section{Квантно рачунарство и различите парадигме}
\label{sec:np}

Квантно рачунарство данас није јединствена архитектура него породица
различитих рачунских парадигми, свака са сопственим физичким темељима,
математичким формализмом и класом проблема за коју је најподесније.
У овом поглављу представљамо две доминантне парадигме које су релевантне
за решавање оптимизационих проблема: \textbf{квантно рачунарство засновано
на колима} и \textbf{квантно жарење}.

% ============================================================
\subsection{Заједничке основе: кубити и суперпозиција}
% ============================================================



Обе парадигме заснивају се на кубиту, квантном аналогу класичног бита.
Стање $n$ кубита описује се вектором у Хилбертовом простору
$\mathcal{H} = (\mathbb{C}^2)^{\otimes n}$ димензије $2^n$:
\[
  |\psi\rangle = \sum_{x \in \{0,1\}^n} \alpha_x |x\rangle, \qquad
  \sum_{x} |\alpha_x|^2 = 1.
\]
Кубит када му се измери вредност може бити само $|0\rangle$ или $|1\rangle$, али пре мерења може бити у суперпозицији свих могућих конфигурација.
На пример кубит у стању $|\psi\rangle = \frac{|0\rangle + |1\rangle}{\sqrt{2}}$ има једнаку вероватноћу да се измери као $|0\rangle$ или $|1\rangle$.
Двокубитни систем у стању $|\psi\rangle = \frac{|00\rangle + |11\rangle}{\sqrt{2}}$ је у суперпозицији стања $|00\rangle$ и $|11\rangle$, али никада неће бити измерен као $|01\rangle$ или $|10\rangle$.
Ова стања одговарају појединачним битовима у класичном рачунарству, али једина разлика је што ми не знамо унапред колика
је вредност сваког бита у стварности све док не измеримо систем.

Суперпозиција омогућује да систем истовремено ``буде'' у свим $2^n$
конфигурацијама пре мерења, управо то чини квантно рачунарство
концептуално различитим од класичног.

% ============================================================
\subsection{Парадигма I: Квантна кола (Gate-Based QC)}
% ============================================================


\subsubsection{Принцип рада}

У моделу квантних кола, рачунање се изводи применом секвенце
\emph{квантних капија},  унитарних трансформација $U$, на почетно
стање $|0\rangle^{\otimes n}$:
\[
  |\psi_{\text{out}}\rangle = U_L \cdots U_2 U_1 \,|0\rangle^{\otimes n}.
\]
Овај израз представља квантно коло са $L$ капија. Које делује на регистар са  $n$ кубита, и примењеује различите трансформације на њега.
Коначни одговор добија се мерењем, које стање пројектује на неки
базни вектор са одређеном вероватноћом.

\subsubsection{Базичне квантне капије}

Свака квантна капија је \emph{унитарна} трансформација ($U U^\dagger = I$),
чиме је рачунање реверзибилно. Матричне репрезентације основних капија:
\[
  H = \frac{1}{\sqrt{2}}\begin{pmatrix}1 & 1\\ 1 & -1\end{pmatrix}, \qquad
  X = \begin{pmatrix}0 & 1\\ 1 & 0\end{pmatrix}, \qquad
  R_Z(\theta) = \begin{pmatrix}e^{-i\theta/2} & 0\\ 0 & e^{i\theta/2}\end{pmatrix}, 
\]

\paragraph{Просте капије (један кубит).}
На слици~\ref{fig:simple-gates} приказано је шест основних
једнокубитних капија. Поред матрица $H$, $X$ и $R_Z$ из горњег
израза, додајемо:
\[
  Y = \begin{pmatrix}0&-i\\i&0\end{pmatrix}, \qquad
  S = \begin{pmatrix}1&0\\0&i\end{pmatrix}, \qquad
  T = \begin{pmatrix}1&0\\0&e^{i\pi/4}\end{pmatrix}, \qquad
  R_Y(\theta) = \begin{pmatrix}\cos\frac{\theta}{2}&-\sin\frac{\theta}{2}\\\sin\frac{\theta}{2}&\cos\frac{\theta}{2}\end{pmatrix}.
\]

\begin{figure}[h]
\centering
\begin{quantikz}
  \lstick{$|q\rangle$} & \gate{H} & \qw & \rstick{\small $\tfrac{|0\rangle\pm|1\rangle}{\sqrt{2}}$}
\end{quantikz}
\quad
\begin{quantikz}
  \lstick{$|q\rangle$} & \gate{X} & \qw & \rstick{\small $|1{\oplus}q\rangle$}
\end{quantikz}
\quad
\begin{quantikz}
  \lstick{$|q\rangle$} & \gate{Y} & \qw & \rstick{\small $iY|q\rangle$}
\end{quantikz}

\vspace{0.6em}

\begin{quantikz}
  \lstick{$|q\rangle$} & \gate{S} & \qw & \rstick{\small $e^{i\pi/2}$ на $|1\rangle$}
\end{quantikz}
\quad
\begin{quantikz}
  \lstick{$|q\rangle$} & \gate{T} & \qw & \rstick{\small $e^{i\pi/4}$ на $|1\rangle$}
\end{quantikz}
\quad
\begin{quantikz}
  \lstick{$|q\rangle$} & \gate{R_Y(\theta)} & \qw & \rstick{\small реална ротација}
\end{quantikz}

\caption{Једнокубитне капије: Хадамар ($H$), Паули-$X$, Паули-$Y$,
         фазна ($S$), $\pi/8$ капија ($T$) и ротација $R_Y(\theta)$.}
\label{fig:simple-gates}
\end{figure}

\paragraph{Двокубитне капије.}
Двокубитне капије делују на пар кубита и могу генерисати заплетеност.
Три кључне:

\begin{figure}[h]
\centering
\begin{quantikz}
  \lstick{$|c\rangle$} & \ctrl{1} & \qw & \rstick{$|c\rangle$} \\
  \lstick{$|t\rangle$} & \targ{}  & \qw & \rstick{$|t\oplus c\rangle$}
\end{quantikz}
\qquad
\begin{quantikz}
  \lstick{$|c\rangle$} & \ctrl{1}   & \qw & \rstick{$(-1)^{ct}|c\rangle$} \\
  \lstick{$|t\rangle$} & \control{} & \qw & \rstick{$(-1)^{ct}|t\rangle$}
\end{quantikz}
\qquad
\begin{quantikz}
  \lstick{$|a\rangle$} & \swap{1} & \qw & \rstick{$|b\rangle$} \\
  \lstick{$|b\rangle$} & \targX{} & \qw & \rstick{$|a\rangle$}
\end{quantikz}
\caption{Двокубитне капије: CNOT (controlled-NOT), CZ (controlled-$Z$)
         и SWAP (размена стања кубита).}
\label{fig:two-qubit-gates}
\end{figure}

\paragraph{Трокубитна Тофолијева капија (CCNOT).}
Тофолијева капија примењује $X$ на циљни кубит само ако су
\emph{оба} управљачка у стању $|1\rangle$.
Будући да реализује NAND операцију, она је функционално потпуна
и доказује да квантна кола могу симулирати произвољно класично рачунање.

\begin{figure}[h]
\centering
\begin{quantikz}
  \lstick{$|a\rangle$} & \ctrl{1} & \qw & \rstick{$|a\rangle$} \\
  \lstick{$|b\rangle$} & \ctrl{1} & \qw & \rstick{$|b\rangle$} \\
  \lstick{$|t\rangle$} & \targ{}  & \qw & \rstick{$|t \oplus (a \wedge b)\rangle$}
\end{quantikz}
\caption{Тофолијева (CCNOT) капија: $X$ се примењује на $|t\rangle$
         само када су $|a\rangle = |b\rangle = |1\rangle$.}
\label{fig:toffoli}
\end{figure}

Оно што можемо да приметимо да нема неке велике разлике између ових капија и класичних логичких капија
на дијаграмима, али је битно да су ове капије реверзибилне, што значи да не губе информацију и да се могу изводити уназад.
Ово је кључна разлика у односу на класичне логичке капије које нису реверзибилне (нпр. AND, OR, NAND). У квантном рачунарству,
реверзибилност је неопходна јер унитарне трансформације морају имати инверз.

Класична AND капија је типичан пример \emph{иреверзибилне} операције.
Она пресликава два улазна бита у један излазни, чиме неповратно губи
информацију. Из излазне вредности $0$ не може се одредити да ли је улаз
био $(0,0)$, $(0,1)$ или $(1,0)$ --- три различите улазне комбинације
дају исти излаз. Таква функција није инјекција и не поседује инверз,
па се не може реализовати унитарном матрицом.

\begin{figure}[h]
\centering
\begin{tikzpicture}[scale=1.0]
  % AND gate body: flat left side, semicircular right side
  \draw[thick] (0,0) -- (0,1.4) -- (0.8,1.4)
        arc[start angle=90, end angle=-90, radius=0.7]
        -- (0,0);
  % Input wires
  \draw[thick] (-1.0,1.0) -- (0,1.0) node[left, at start] {$A$};
  \draw[thick] (-1.0,0.4) -- (0,0.4) node[left, at start] {$B$};
  % Output wire
  \draw[thick] (1.5,0.7) -- (2.3,0.7) node[right] {$A \wedge B$};
  % Gate label
  \node at (0.55,0.7) {\small AND};
\end{tikzpicture}
\qquad\qquad
% Truth table
\begin{tabular}[b]{cc|c}
  \toprule
  $A$ & $B$ & $A \wedge B$ \\
  \midrule
  0 & 0 & 0 \\
  0 & 1 & 0 \\
  1 & 0 & 0 \\
  1 & 1 & 1 \\
  \bottomrule
\end{tabular}
\caption{AND капија и њена таблица истинитости. Три улазне комбинације
  $(0,0)$, $(0,1)$ и $(1,0)$ дају исти излаз $0$, па функција
  није инјекција, информација о улазу се губи и операција
  није реверзибилна.}
\label{fig:and-gate}
\end{figure}

\subsubsection{Хардверске платформе}

Неколико компанија данас нуди комерцијално доступне квантне рачунаре
засноване на колима:

\textbf{IBM Quantum} користи суперпроводљиве кубите охлађене близу
апсолутне нуле (${\sim}15\,\mathrm{mK}$). Кроз платформу IBM Quantum
Experience корисници могу покретати програме написане у Qiskit-у
директно на стварном хардверу. Серија Eagle/Heron процесора достигла
је преко 1000 физичких кубита (Condor, 2023), мада је број грешака
и даље ограничавајући фактор.

\textbf{Google Quantum AI} такође гради суперпроводљиве процесоре.
Чип Sycamore (54 кубита, 2019) коришћен је у демонстрацији
,,квантне надмоћи'', задатку за који је тврђено да класичном
рачунару треба 10.000 година, а процесор га је решио за 200 секунди
(тврдња је касније оспорена у литератури).
Новији чип Willow (2024) показује напредак у корекцији грешака.

\textbf{IonQ} ради на бази заробљених јона (trapped-ion технологија):
атоми итербијума одржавају се у вакуумској замци електромагнетним
пољем, а капије се примењују ласерским пулсовима. Иако су тренутно
спорији од суперпроводљивих система, кохеренције јони имају знатно
нижу стопу грешака и дуже вријеме кохеренце.

\textbf{Quantinuum} (спој Honeywell Quantum Solutions и Cambridge
Quantum) такође примењује приступ заробљених јона и истиче се међу
највишим вредностима квантног волумена --- мере која узима у обзир
број кубита, повезаност и стопу грешака истовремено.

\textbf{Остали актери.} Компаније попут Rigetti (суперпроводљивост),
PsiQuantum (фотонски кубити) и Intel (silicon spin qubits) истражују
алтернативне физичке имплементације у потрази за скалабилнијим решењем.

\subsubsection{Ограничења у NISQ ери}


Тренутни квантни рачунари са колима налазе се у такозваној
\textit{NISQ ери}~\cite{preskill2018} (Noisy Intermediate-Scale Quantum):
имају 50--1000 кубита, али су подложни грешкама и немају потпуну
корекцију грешака. Дубина кола ограничена је временом декохеренције.

% ============================================================
\subsection{Парадигма II: Квантно жарење (Quantum Annealing)}
% ============================================================

% --- Опис ---
% Квантно жарење је специјализована парадигма за оптимизационе
% проблеме. Инспирисана је класичним симулираним жарењем, али
% користи квантно тунелирање уместо термалних флуктуација.

\subsubsection{Класично симулирано жарење}

Симулирано жарење (Kirkpatrick et al., 1983) је метахеуристички
алгоритам инспирисан процесом физичког жарења метала: загревање
уклања дефекте у кристалној решетки, а споро хлађење омогућује
систему да нађе нискоенергетско стање.

\textbf{Идеја алгоритма.} Циљ је минимизовати функцију цене $E(s)$
над дискретним простором стања $s$. У свакој итерацији алгоритам:
\begin{enumerate}[leftmargin=*, label=\arabic*.]
  \item предложи мали случајни помак у суседно стање $s'$;
  \item ако је $E(s') < E(s)$, прихвати помак безусловно;
  \item ако је $E(s') \geq E(s)$, прихвати помак са вероватноћом
        $P = e^{-\Delta E / T}$, где је $\Delta E = E(s') - E(s)$
        и $T$ тренутна ``температура''.
\end{enumerate}
Температура се постепено смањује по \emph{распореду хлађења}
(нпр. $T_k = T_0 / \log(1+k)$). Висока температура на почетку
омогућује велике скокове преко лоших локалних решења; ниска
температура на крају усмерава претрагу ка финалном минимуму.

\textbf{Ограничење.} Класично жарење прескаче баријере искључиво
захваљујући температурним флуктуацијама. Уска, висока баријера
захтева велику температуру --- а висока температура истовремено
деградира прецизност финалног решења. Овај компромис отежава
проналажење глобалног минимума у пејзажима са бројним уским
локалним минимумима.

\subsubsection{Принцип рада}

Идеја квантног жарења потиче из аналогије са класичним
\emph{симулираним жарењем} (Kirkpatrick et al., 1983), али уместо
термалних флуктуација користи \emph{квантно тунеловање} — способност
квантног система да пролази кроз енергетске баријере уместо да их
прескаче. Ово је кључна предност над класичним жарењем када
енергетски пејзаж има бројне уске, дубоке локалне минимуме.

\textbf{Почетни Хамилтонијан} $\hat{H}_0$ бира се тако да му је
основно стање лако за припремити. Стандардан избор је попречно
магнетно поље:
\[
  \hat{H}_0 = -\sum_{i=1}^n \hat{X}_i,
\]
чије је основно стање равномерна суперпозиција свих $2^n$ стања,
$|+\rangle^{\otimes n}$.

\textbf{Проблемски Хамилтонијан} $\hat{H}_P$ кодира функцију коју
оптимизујемо: његово основно стање одговара оптималном решењу.
За QUBO/Изинг проблеме:
\[
  \hat{H}_P = \sum_{i} h_i \hat{Z}_i + \sum_{i < j} J_{ij} \hat{Z}_i \hat{Z}_j.
\]

\textbf{Адијабатска еволуција.} Систем се спором линеарном
интерполацијом води из $\hat{H}_0$ у $\hat{H}_P$:
\[
  \hat{H}(t) = \left(1 - \frac{t}{T}\right)\hat{H}_0
             + \frac{t}{T}\,\hat{H}_P, \qquad t \in [0, T].
\]
У тренутку $t = 0$ систем је у основном стању $\hat{H}_0$.
На крају ($t = T$) мерење даје стање блиско основном стању $\hat{H}_P$,
тј.\ (приближно) оптималном решењу.

\begin{theorem}[Адијабатска теорема]
Нека је $\Delta(t) = E_1(t) - E_0(t) > 0$ спектрални јаз
(разлика енергија основног и првог побуђеног стања у тренутку $t$).
Ако је укупно време жарења $T$ довољно велико у односу на
\[
  T \;\gg\; \frac{\max_t \|\partial_t \hat{H}(t)\|}{\min_t \Delta(t)^2},
\]
систем иницијализован у основном стању $\hat{H}_0$ остаје у
тренутном основном стању током целог процеса жарења.
\end{theorem}

\begin{remark}
Кључни изазов адијабатског рачунарства јесте \emph{спектрални јаз}:
ако $\Delta(t)$ постане веома мали током еволуције, потребно је
експоненцијално велико $T$, чиме се предност губи.
За НП-тешке проблеме очекује се да минимални јаз опада
полиномски или чак експоненцијално са величином инстанце.
\end{remark}

\textbf{Квантно тунелирање.} За разлику од класичног жарења које
„прескаче" преко баријере захваљујући температури, квантни систем
тунелира \emph{кроз} баријеру захваљујући суперпозицији.
Уска, висока баријера коју класично жарење тешко прелази може
бити прозирна за квантно тунелирање, па квантно жарење може
ефикасније наћи глобални минимум у одређеним пејзажима.

На слици~\ref{fig:annealing-schedule} приказан је типичан
распоред жарења: попречно поље $\Gamma(t)$ постепено слаби,
а проблемски члан $\lambda(t)$ јача.

\begin{figure}[h]
\centering
\begin{tikzpicture}[scale=1.1]
  \draw[->] (0,0) -- (5.2,0) node[right] {$t/T$};
  \draw[->] (0,0) -- (0,2.4) node[above] {амплитуда};
  % Gamma(t) = 1 - t/T (poprecno polje opada)
  \draw[blue!70!black, thick] (0,2) -- (5,0)
        node[right, blue!70!black] {$\Gamma(t) = 1 - t/T$};
  % lambda(t) = t/T (problemski clan raste)
  \draw[red!70!black, thick, dashed] (0,0) -- (5,2)
        node[right, red!70!black] {$\lambda(t) = t/T$};
  % ose
  \foreach \x/\l in {0/{$0$}, 2.5/{$T/2$}, 5/{$T$}}
    \draw (\x,0.05) -- (\x,-0.05) node[below] {\l};
  \foreach \y in {1,2}
    \draw (0.05,\y) -- (-0.05,\y) node[left] {$\y$};
\end{tikzpicture}
\caption{Распоред жарења: попречно поље $\Gamma(t)$ (квантне флуктуације)
         постепено слаби, а проблемски члан $\lambda(t)$ јача.
         На крају ($t = T$) доминира $\hat{H}_P$.}
\label{fig:annealing-schedule}
\end{figure}

\subsubsection{QUBO формулација}

% --- QUBO ---
% Квантно жарење природно решава QUBO проблеме.
% QUBO: minimize x^T Q x, x ∈ {0,1}^n.
% Исинговом трансформацијом (z = 2x-1) добија се Изингов модел
% директно реализован у хардверу.

Природни улазни формат за квантне жараче је
\textbf{QUBO} (Quadratic Unconstrained Binary Optimization):
\[
  \min_{\mathbf{x} \in \{0,1\}^n} \;\mathbf{x}^\top Q\, \mathbf{x},
  \qquad Q \in \mathbb{R}^{n \times n}.
\]

Супституцијом $z_i = 2x_i - 1 \in \{-1,+1\}$ добија се еквивалентни
\emph{Изингов модел}, директно реализован у хардверу квантног жарача:
\[
  \hat{H}_P = \sum_{i} h_i \hat{Z}_i + \sum_{i < j} J_{ij} \hat{Z}_i \hat{Z}_j.
\]


\subsubsection{Хардвер: D-Wave}

D-Wave Systems је тренутно једини комерцијални произвођач квантних
жарача. Њихов процесор \textbf{Advantage} (2020) садржи преко 5000
суперпроводљивих кубита охлађених на ${\sim}15\,\mathrm{mK}$, а
генерација \textbf{Advantage2} (у развоју) циља преко 7000 кубита са
побољшаним графом повезаности.

\textbf{Архитектура повезаности.}
Кубити нису потпуно повезани --- хардверски граф је реалан и редак
(sparse). Процесор Advantage користи \emph{Pegasus} граф у коме је
сваки кубит директно повезан са највише 15 суседа, за разлику од
старије \emph{Chimera} архитектуре (максимум 6 веза по кубиту).

\textbf{Уграђивање (Embedding).}
Пошто произвољан проблемски граф обично није подграф Pegasus-а,
неопходно је \emph{уграђивање}: логички кубит проблема представља
се ланцем физичких кубита везаних јаком феромагнетном везом.
Уграђивање може значајно повећати стварни број потребних физичких
кубита и представља практично ограничење за велике инстанце.

\textbf{Облачни приступ и примене.}
D-Wave нуди приступ преко \emph{Leap} облачне платформе са Python
библиотеком \texttt{dwave-ocean-sdk}. Међу познатим применама истичу
се: оптимизација рутирања возила (Volkswagen), финансијско
портфолио управљање и проблеми распоређивања у логистици.

% ============================================================
\subsection{Поређење парадигми}
% ============================================================

% TODO: Табела поређења:
% | Особина          | Кола (Gate-Based) | Жарење (Annealing) |
% |------------------|-------------------|--------------------|
% | Модел            | Унитарна кола     | Адијабатска еволуција|
% | Улаз             | Квантни програм   | QUBO/Изинг матрица |
% | Излаз            | Мерење на крају   | Основно стање      |
% | Платформе        | IBM, Google, IonQ | D-Wave             |
% | Намена           | Универзалан       | Оптимизација       |
% | Кубити (2024)    | ~1000 (физички)   | >5000              |
% | Корекција грешака| У развоју         | Не постоји         |

\begin{table}[ht]
  \centering
  \caption{Поређење две главне квантне парадигме.}
  \label{tab:paradigme}
  \renewcommand{\arraystretch}{1.3}
  \begin{tabular}{lll}
    \toprule
    \textbf{Особина}         & \textbf{Кола (Gate-Based)} & \textbf{Жарење (Annealing)} \\
    \midrule
    Рачунски модел           & Унитарна кола              & Адијабатска еволуција \\
    Улазни формат            & Квантни програм            & QUBO / Изингов модел \\
    Платформе                & IBM, Google, IonQ          & D-Wave \\
    Примарна намена          & Универзалан                & Комбинаторна оптимизација \\
    Број кубита (2024)       & ${\sim}1000$ физичких      & ${>}5000$ \\
    Корекција грешака        & У развоју                  & Не постоји \\
    \bottomrule
  \end{tabular}
\end{table}

Иако се парадигме разликују у механизму, оне нису потпуно раздвојене:
QAOA алгоритам са колима, инспирисан је
управо адијабатским жарењем и решава исте QUBO проблеме.

% ============================================================
%  Поглавље 3: Проблем MaxCut
% ============================================================
\section{Проблем MaxCut}
\label{sec:maxcut}

% --- Формална дефиниција ---
% Дефинишите граф G=(V,E) са тежинама. Дефинишите рез (бисекцију)
% скупа V у два скупа S и V\S. Циљ: максимизовати тежину
% ивица које прелазе рез.

\subsection{Формална дефиниција}

\begin{definition}[MaxCut]
Нека је $G = (V, E, w)$ неусмерени граф са тежинама $w : E \to \mathbb{R}_{+}$.
\emph{Рез} је партиција $V = S \cup \bar{S}$, $S \cap \bar{S} = \emptyset$.
Проблем \textsc{MaxCut} тражи партицију која максимизује
\[
  \text{cut}(S) = \sum_{\{u,v\} \in E,\; u \in S,\; v \in \bar{S}} w(u,v).
\]
\end{definition}

% --- Пример на малом графу ---
% Нацртајте граф са 5--6 чворова (TikZ), прикажите оптималан рез
% и вредност циља. Ово ће послужити и у каснијим поглављима
% (Z3 кодирање, QAOA).

\subsection{Илустративни пример}

% TODO: Нацртати TikZ граф, означити оптималан рез испрекиданом линијом.
% Препоручени пример: K4 граф (потпуни граф на 4 чвора), MaxCut = 4.

\subsection{НП-тешкоћа MaxCut-а}

% --- Редукција из 3-SAT (или NAE-3SAT) ---
% Скицирајте доказ НП-тешкоћа MaxCut-а полиномском редукцијом
% из НП-потпуног проблема (нпр. NAE-3SAT → MaxCut).
% Укључите референцу на Карп 1972.

\begin{theorem}[Карп, 1972~\cite{karp1972}]
Проблем \textsc{MaxCut} је \textsf{NP}-тешки.
\end{theorem}

% TODO: Скица доказа редукцијом.

\subsection{Класични апроксимациони алгоритми}

% --- Гоеманс-Вилијамсон алгоритам ---
% Опишите SDP релаксацију и хиперраванску рандомизацију.
% Наведите гаранцију апроксимације: ≥ 0.878 · OPT.
% Ово мотивише питање: може ли квантни приступ да уради боље?

Гоеманс и Вилијамсон~\cite{goemans1995} предложили су алгоритам заснован
на семидефинитном програмирању (SDP) са гаранцијом апроксимације:
\[
  \text{cut}(S_{\text{alg}}) \;\geq\; 0.878 \cdot \text{OPT}.
\]

% TODO: Опис SDP релаксације и хиперраванске рандомизације.
% Напомена: под условом P ≠ NP и UGC хипотезе, ово је оптимална
% апроксимација за полиномске алгоритме.

% ============================================================
%  Поглавље 4: Класични приступ — Z3 решавач
% ============================================================
\section{Класични приступ: Z3 решавач и SMT}
\label{sec:z3}

% --- Увод у SAT/SMT решаваче ---
% Кратко представите SAT решаваче (DPLL, CDCL), а затим SMT
% (Satisfiability Modulo Theories) као проширење са додатним теоријама
% (аритметика, низови, битвектори...).

\subsection{SAT и SMT решавачи}

% TODO: Опишите DPLL/CDCL архитектуру у 1--2 параграфа.
% Поменути Nelson-Oppen оквир за комбиновање теорија у SMT.

\subsection{Z3: Архитектура и могућности}

% --- Опишите Z3 ---
% Z3 је Microsoft Research SMT решавач (de Moura & Bjørner, 2008).
% Подржава: линеарну аритметику, низове, битвекторе, квантификаторе...
% Сучеља: Python API (z3-solver), C, .NET.

% TODO: Укључити кратак блок Python кода са z3-solver:
% deklaracija varijabli, dodavanje ograničenja, poziv solve().

\subsection{Кодирање MaxCut-а у Z3}

% --- Кодирање ---
% За граф G=(V,E), уведите булове променљиве x_v ∈ {0,1} за сваки чвор.
% Ивица {u,v} прелази рез ако x_u ≠ x_v, тј. x_u XOR x_v = 1.
% Циљ: максимизовати Σ_{(u,v)∈E} (x_u XOR x_v).
%
% SMT кодирање: Int или Bool за x_v, Optimize() за максимизацију.

\begin{example}[MaxCut на $K_4$ у Z3 (Python)]
\begin{verbatim}
from z3 import *
x = [Int(f'x{i}') for i in range(4)]
opt = Optimize()
for v in x:
    opt.add(Or(v == 0, v == 1))
edges = [(0,1),(0,2),(0,3),(1,2),(1,3),(2,3)]
cut = Sum([If(x[u] != x[v], 1, 0) for u,v in edges])
opt.maximize(cut)
print(opt.check(), opt.model())
\end{verbatim}
\end{example}

% TODO: Показати излаз за K4 (оптимум = 4) и прокоментарисати.

\subsection{Ограничења SMT приступа}

% --- Скалабилност ---
% Употреба Z3 за MaxCut је прецизна (даје оптимум), али NP-тешкоћа
% проблема значи да се време извршавања у најгорем случају
% експоненцијално повећава са бројем чворова.
% Приложити табелу: n=10,20,50 чворова vs. CPU време у Z3.

% TODO: Мерења перформанси Z3 за случајне графове растуће величине.
% Ово мотивише прелаз на квантни приступ.

\begin{remark}
Z3 је погодан алат за верификацију малих инстанци и за генерисање
референтних решења, али не може заменити апроксимационе или квантне
алгоритме за велике улазе.
\end{remark}

% ============================================================
%  Поглавље 5: Основе квантног рачунарства
% ============================================================
\section{Основе квантног рачунарства}
\label{sec:kvant}

% --- Циљ поглавља ---
% Дати минималну али самодовољну основу квантног рачунарства
% неопходну за разумевање QAOA у следећем поглављу.
% Избегавати претерану физичку формалност; фокус на рачунарски аспект.

\subsection{Кубити и суперпозиција}

% --- Кубит ---
% Кубит је квантна аналогија класичног бита. Стање кубита је вектор
% у дводимензионалном комплексном Хилбертовом простору:
%   |ψ⟩ = α|0⟩ + β|1⟩,  |α|² + |β|² = 1.
% Мерење даје 0 са вероватноћом |α|² и 1 са вероватноћом |β|².

Стање кубита описује се вектором у Хилбертовом простору $\mathcal{H} = \mathbb{C}^2$:
\[
  |\psi\rangle = \alpha|0\rangle + \beta|1\rangle, \quad
  \alpha, \beta \in \mathbb{C}, \quad |\alpha|^2 + |\beta|^2 = 1.
\]

% TODO: Илустровати Блохову сферу (TikZ или слика).
% Поменути да систем од n кубита живи у простору 2ⁿ димензија.

\subsection{Квантна заплетеност}

% --- Entanglement ---
% Два кубита су заплетена ако се стање система не може записати
% као тензорски производ стања појединачних кубита.
% Пример: Белово стање (|00⟩ + |11⟩)/√2.
% Заплетеност омогућује корелације без класичне комуникације.

\begin{example}[Белово стање]
\[
  |\Phi^+\rangle = \frac{|00\rangle + |11\rangle}{\sqrt{2}}
\]
Мерење prvog кубита тренутно одређује стање другог, без обзира на растојање.
\end{example}

\subsection{Квантна кола и капије}

% --- Квантне капије ---
% Квантне капије су унитарне трансформације. Набројите:
% H (Хадамар), X (NOT), CNOT, Rz(θ), CZ.
% Нагласите да је Хадамар капија кључна за стварање суперпозиције.

\[
  H = \frac{1}{\sqrt{2}}\begin{pmatrix}1 & 1\\ 1 & -1\end{pmatrix}, \qquad
  X = \begin{pmatrix}0 & 1\\ 1 & 0\end{pmatrix}, \qquad
  R_Z(\theta) = \begin{pmatrix}e^{-i\theta/2} & 0\\ 0 & e^{i\theta/2}\end{pmatrix}.
\]

% TODO: Приказати једноставно квантно коло (нпр. Bell state) помоћу
% qcircuit пакета или TikZ.

\subsection{Квантна предност и NISQ ера}

% --- NISQ ---
% Прескил 2018: тренутни квантни рачунари имају 50--1000 зашумљених
% кубита (Noisy Intermediate-Scale Quantum, NISQ).
% Нису довољно велики за потпуну квантну корекцију грешака,
% али омогућују хибридне квантно-класичне алгоритме.

% TODO: Поменути тренутне хардверске платформе:
% IBM Quantum, Google Sycamore, IonQ, Rigetti.
% Навести тренутне рекорде у броју кубита и дубини кола.

\begin{remark}
Квантна предност за НП-тешке проблеме теоријски није доказана.
Истраживања показују да NISQ алгоритми могу пружити практичну предност
на специфичним инстанцама пре него теоријску асимптотску супериорност.
\end{remark}

% ============================================================
%  Поглавље 6: QAOA — Квантни алгоритам за оптимизацију
% ============================================================
\section{QAOA и MaxCut}
\label{sec:qaoa}

QAOA (\emph{Quantum Approximate Optimization Algorithm}) предложили су
Farhi, Goldstone и Gutmann као хибридни квантно-класични алгоритам за
комбинаторну оптимизацију~\cite{farhi2014}. За наше потребе није
пресудан детаљан профил перформанси QAOA, већ чињеница да он показује
како се \emph{иста} проблемска формулација, добијена у облику
QUBO/Изинговог модела, може користити и у моделу квантних кола.

\textbf{Веза QUBO $\to$ Изинг $\to$ QAOA.}
Сваки QUBO проблем
се супституцијом $x_i = (1-s_i)/2$ директно пресликава у Изингов
Хамилтонијан облика
$H = \sum_i h_i s_i + \sum_{i<j} J_{ij} s_i s_j$,
где спинови $s_i \in \{-1,+1\}$ одговарају кубитима у стањима
$|0\rangle$ и $|1\rangle$.
QAOA затим узима тај Хамилтонијан \emph{директно} као проблемски
оператор $\hat{H}_C$ без икакве додатне трансформације.
Другим речима, QUBO матрица $Q$ потпуно одређује квантно коло QAOA-е:
дијагонални елементи $Q_{ii}$ дају локална поља $h_i$,
а вандијагонални $Q_{ij}$ дају спрезања $J_{ij}$.

\subsection{Од MaxCut-а до проблемског Хамилтонијана}

За MaxCut уводимо спин варијабле $z_v \in \{-1,+1\}$ за сваки чвор
$v \in V$. Ивица $\{u,v\}$ прелази рез ако и само ако $z_u z_v = -1$,
па је величина реза
\[
  C(\mathbf{z}) = \sum_{\{u,v\}\in E} \frac{1-z_u z_v}{2}.
\]
Ова циљна функција директно прелази у проблемски Хамилтонијан
\[
  \hat{H}_C = \frac{1}{2}\sum_{\{u,v\}\in E}\bigl(\hat{I}-\hat{Z}_u\hat{Z}_v\bigr).
\]

Кључна ствар је да је то \emph{иста} енергетска структура коју смо већ
добили из QUBO формулације у поглављу~\ref{sec:qubo}. Зато QAOA није
нова формулација проблема, већ нови начин да се оптимизује већ познати
Хамилтонијан.


\subsection{Зашто је QAOA важан}

У контексту наше дискусије, QAOA је важан из три разлога.
\begin{itemize}[nosep]
  \item Показује да се MaxCut може решавати у моделу квантних кола над
        истим QUBO/Изинговим описом који користе квантни анилери.
  \item Илуструје хибридни образац: квантни део претражује простор
        стања, а класични део управља оптимизацијом.
  \item Пружа природно место за поређење са Z3: QAOA тражи добро
        решење варијационо, док Z3 може егзактно да провери мале
        инстанце и исправност добијеног bitstring-а.
\end{itemize}

Практично ограничење QAOA у NISQ окружењу јесте то што дубина кола,
шум и цена класичне оптимизационе петље брзо постају доминантни.
Зато је у овом раду QAOA пре свега \emph{концептуални мост} између
QUBO формулације и квантног рачунарства заснованог на колима и указујемо на њега 
као пример, како се исти проблем може решавати у различитим моделима рачунања.

% ============================================================
%  Поглавље 7: Веза са аутоматским резоновањем
% ============================================================
\section{Квантно рачунарство и аутоматско резоновање}
\label{sec:ar}

% --- Циљ поглавља ---
% Повезати квантне технике са темом семинара: аутоматским резоновањем.
% Показати да постоје дубљи концептуални мостови, а не само
% аналогија по примени.

\subsection{SAT и квантно решавање}

Гроверов алгоритам претражује неструктурисани простор $N$ стања
у $O(\sqrt{N})$ корака, наспрам класичних $O(N)$~\cite{grover1997}.
За SAT са $n$ варијабли то даје убрзање:
\[
  O(2^n) \;\longrightarrow\; O(2^{n/2}).
\]

Интуиција је следећа. Булова формула
\[
  \varphi(x_1,\dots,x_n)
\]
посматра се као функција која за сваку доделу
$\mathbf{x}\in\{0,1\}^n$ враћа $1$ ако је формула задовољена,
а $0$ иначе. Тада се конструише \emph{квантни оракл}
\[
  O_\varphi \ket{\mathbf{x}} = (-1)^{\varphi(\mathbf{x})}\ket{\mathbf{x}},
\]
који не мења битски низ, већ мења \emph{фазу} оних стања која
представљају задовољавајуће доделе. Након тога Гроверова
итерација појачава амплитуде „добрих" стања и потискује остала,
па мерење са великом вероватноћом враћа решење.

Ако постоји $M$ задовољавајућих додела, сложеност је
$O(\sqrt{2^n/M})$, па је најчешће цитирани случај
$M=1$ или мало $M$, где добијамо $O(2^{n/2})$.
Међутим, ово \emph{не} значи да SAT постаје лак у полиномском смислу:
убрзање је квадратно, али је и даље експоненцијално, што је у складу
са релативизованим доњим границама за неструктурисану претрагу~\cite{bennett1997}.


\subsection{Хибридни квантно-класични приступи}

% --- Хибридни приступи ---
% Модеран приступ: QAOA и VQE (Variational Quantum Eigensolver) нису
% чисто квантни --- они користе класичне оптимизаторе за подешавање
% квантних параметара. Ово је директна аналогија са SMT решавачима
% које воде класични CDCL алгоритми.

% TODO: Нацртати дијаграм хибридне петље:
%   Квантно коло → Мерење → Класична оптимизација → (итерација)
% Поредити са CDCL петљом у SAT решавачима.

\subsection{SMT теорије и квантни Хамилтонијани}

Свака SMT формула може се посматрати као скуп ограничења над
променљивама, док се QUBO/Изингова формулација посматра као
енергетска функција чији минимум репрезентује решење~\cite{barrett2009,lucas2014,glover2022qubo}.
Та паралела је суштинска:
\begin{itemize}[nosep]
  \item SMT ограничења $C_i(\mathbf{x})$ одговарају penalty члановима
        $P_i(\mathbf{x})$ у QUBO Хамилтонијану,
  \item модел SMT формуле одговара основном стању
        (стању минималне енергије),
  \item теоријски решавач у SMT-у игра улогу „симболичког физичара":
        формално проверава да ли је кодирање коректно пре него што
        га предамо квантном или хеуристичком оптимизатору.
\end{itemize}

У том смислу, квантно жарење не замењује формалне методе, већ их
допуњује: квантни уређај претражује простор стања, док SMT решавач
проверава да ли енергетска функција заиста репрезентује оригинални
комбинаторни проблем.

\begin{figure}[h]
\centering
\begin{tikzpicture}[
  node distance=1.4cm and 0.9cm,
  box/.style={
    draw,
    rounded corners=4pt,
    thick,
    align=center,
    text width=2.45cm,
    minimum width=2.45cm,
    minimum height=1.0cm,
    font=\small,
    fill=gray!8
  },
  classical/.style={box, fill=blue!8, draw=blue!50!black},
  qubo/.style={box, fill=orange!12, draw=orange!70!black},
  quantum/.style={box, fill=red!8, draw=red!65!black},
  verify/.style={box, fill=green!10, draw=green!45!black},
  flow/.style={->, thick}
]
  \node[classical] (orig) {Полазни\\проблем};
  \node[classical, right=of orig] (smt) {SMT\\модел};
  \node[qubo, right=of smt] (qubo) {QUBO /\\Изинг};
  \node[quantum, right=of qubo] (solver) {Квантни\\решавач};
  \node[verify, below=1.4cm of qubo] (check) {Z3: провера\\решења};

  \draw[flow] (orig) -- (smt);
  \draw[flow] (smt) -- (qubo);
  \draw[flow] (qubo) -- (solver);
  \draw[flow] (orig.south) |- (check.west);
  \draw[flow] (solver.south) |- node[pos=0.3, right, font=\scriptsize] {$\hat{\mathbf{x}}$} (check.east);
\end{tikzpicture}
\caption{Хибридни ток рада у коме се полазни ограничени проблем најпре
преводи у SMT и QUBO/Изингову форму, затим решава квантно или
хеуристички, а добијени кандидат се на крају проверава помоћу Z3.}
\label{fig:qubo-z3-pipeline}
\end{figure}

Слика~\ref{fig:qubo-z3-pipeline} илуструје централну тезу овог рада:
QUBO није замена за аутоматско резоновање, већ интерфејс између
симболичког описа проблема и квантне оптимизације, док Z3 остаје
формални механизам који затвара петљу и проверава коректност добијених
решења.

\subsection{Z3 као верификатор QUBO ограничења}

Кључна слабост QUBO формулације је у томе што су ограничења
уграђена \emph{индиректно}, преко penalty чланова~\cite{lucas2014,glover2022qubo}.
Ако је знак погрешан, ако недостаје један квадратни члан, или ако је
коефицијент казне сувише мали, минимум енергије више не мора да
одговара важећем решењу оригиналног проблема.
Z3 је зато користан не само као решавач, већ и као
\emph{верификатор исправности саме QUBO трансформације}~\cite{demoura2008,bjorner2014}.

Нека је оригинални проблем дат у облику
\[
  \min g(\mathbf{x})
  \qquad \text{уз услове} \qquad
  C_1(\mathbf{x}), \dots, C_m(\mathbf{x}),
\]
а QUBO апроксимација као
\[
  f(\mathbf{x}) = g(\mathbf{x}) + \sum_{j=1}^m \lambda_j P_j(\mathbf{x}),
  \qquad \mathbf{x} \in \{0,1\}^n.
\]
Тада Z3 може да послужи у најмање три различита корака верификације.

\textbf{1. Провера семантике penalty члана.}
За свако ограничење треба доказати да penalty терм заиста
разликује дозвољена и недозвољена стања:
\[
  C_j(\mathbf{x}) \iff P_j(\mathbf{x}) = 0,
  \qquad
  \neg C_j(\mathbf{x}) \Rightarrow P_j(\mathbf{x}) > 0.
\]
На пример, за ограничење „тачно један" $x_i + x_j = 1$,
penalty члан
\[
  P(\mathbf{x}) = \lambda\,(1 - x_i - x_j + 2x_ix_j)
\]
је исправан ако Z3 покаже да је формула
\[
  x_i,x_j \in \{0,1\}
  \;\land\;
  \bigl((x_i+x_j=1 \land P \neq 0)
  \lor
  (x_i+x_j\neq 1 \land P = 0)\bigr)
\]
незадовољива. На тај начин добијамо \emph{контрапример ако кодирање
није добро}, уместо да се ослањамо само на ручни алгебарски развој.

\begin{example}[Провера penalty члана помоћу Z3]
\begin{lstlisting}[style=python, caption={Z3 provera da penalty clan korektno kodira ogranicenje $x_i + x_j = 1$.}]
from z3 import *

xi, xj, lam = Ints('xi xj lam')

# Penalty term for "exactly one"
P = lam * (1 - xi - xj + 2 * xi * xj)

s = Solver()
s.add(Or(xi == 0, xi == 1))
s.add(Or(xj == 0, xj == 1))
s.add(lam > 0)

# Search for a counterexample:
# valid assignment with nonzero penalty, or
# invalid assignment with nonpositive penalty
s.add(Or(
    And(xi + xj == 1, P != 0),
    And(xi + xj != 1, P <= 0)
))

print(s.check())  # unsat
\end{lstlisting}
\end{example}

Резултат \texttt{unsat} значи да контрапример не постоји:
penalty члан је нула тачно на дозвољеним стањима, а строго позитиван
на недозвољеним.

\textbf{2. Провера да је казна довољно велика.}
Чак и ако је $P_j(\mathbf{x}) = 0$ тачно на дозвољеним стањима,
коефицијент $\lambda_j$ мора бити довољно велик да ниједно
недозвољено стање не постане енергетски повољније од исправног решења~\cite{lucas2014,glover2022qubo}.
За мале и средње инстанце Z3 може прво да нађе оптимално дозвољено
решење $\mathbf{x}^\star$, а затим да провери да ли постоји
контрапример облика
\[
  \exists\, \mathbf{x}\in\{0,1\}^n :
  \neg C(\mathbf{x}) \land f(\mathbf{x}) \leq f(\mathbf{x}^\star).
\]
Ако је ова формула \texttt{unsat}, онда ниједно кршење ограничења
не може да „пробије" легално решење у QUBO енергији.
То је формални доказ да је изабрана penalty скала довољна.

\begin{example}[Провера да је penalty коефицијент довољно велик]
\begin{lstlisting}[style=python, caption={Z3 provera da penalizacija eliminise nedozvoljena resenja.}]
from z3 import *

# Original constrained problem:
# maximize 3*x0 + 2*x1  subject to x0 + x1 <= 1
x0, x1 = Ints('x0 x1')

opt = Optimize()
opt.add(Or(x0 == 0, x0 == 1), Or(x1 == 0, x1 == 1))
opt.add(x0 + x1 <= 1)

reward = 3 * x0 + 2 * x1
h = opt.maximize(reward)

assert opt.check() == sat
m = opt.model()
best_reward = m.eval(reward).as_long()   # 3
best_energy = -best_reward               # -3

def violates_feasible_optimum(lam_value):
    y0, y1 = Ints('y0 y1')
    s = Solver()
    s.add(Or(y0 == 0, y0 == 1), Or(y1 == 0, y1 == 1))

    qubo = -(3 * y0 + 2 * y1) + lam_value * y0 * y1

    # Look for an infeasible assignment that is at least as good
    # as the best feasible solution
    s.add(y0 + y1 > 1)
    s.add(qubo <= best_energy)

    return s.check()

print(violates_feasible_optimum(2))  # sat
print(violates_feasible_optimum(3))  # unsat
\end{lstlisting}
\end{example}

Овде \texttt{sat} за $\lambda = 2$ значи да и даље постоји
недозвољено стање које достиже енергију најбољег легалног решења,
док \texttt{unsat} за $\lambda = 3$ формално потврђује да је
та казна довољна.

\textbf{3. Верификација резултата квантног решавача.}
Када QAOA или квантни анилер врате кандидатно решење
$\hat{\mathbf{x}}$~\cite{farhi2014,preskill2018}, Z3 може независно да провери:
\begin{itemize}[nosep]
  \item да ли $\hat{\mathbf{x}}$ задовољава оригинална ограничења,
  \item колика је вредност оригиналне циљне функције $g(\hat{\mathbf{x}})$,
  \item да ли постоји боље легално решење за исту малу инстанцу.
\end{itemize}
Тако Z3 постаје \emph{референтни оракл} за процену квалитета
квантних и хеуристичких метода: квантни алгоритам тражи ниску
енергију, а SMT верификатор потврђује да та ниска енергија има
жељено семантичко значење.

\begin{example}[Верификација кандидатног bitstring-а]
\begin{lstlisting}[style=python, caption={Z3 provera kandidatnog resenja dobijenog kvantnim ili heuristickim resavacem.}]
from z3 import *

# Candidate returned by an annealer or QAOA
x_hat = [0, 1, 0]
weights = [5, 4, 1]

# Original problem: maximize 5*x0 + 4*x1 + x2
# subject to exactly one variable being 1
x = [Int(f'x{i}') for i in range(3)]

# 1) Check whether x_hat is feasible
s = Solver()
for v in x:
    s.add(Or(v == 0, v == 1))
for i, bit in enumerate(x_hat):
    s.add(x[i] == bit)
s.add(Sum(x) == 1)

print(s.check())  # sat

# 2) Evaluate the returned solution
hat_value = sum(weights[i] * x_hat[i] for i in range(3))
print(hat_value)  # 4

# 3) Search for a better feasible solution
opt = Optimize()
for v in x:
    opt.add(Or(v == 0, v == 1))
opt.add(Sum(x) == 1)

objective = Sum([weights[i] * x[i] for i in range(3)])
h = opt.maximize(objective)

if opt.check() == sat:
    m = opt.model()
    best_x = [m.eval(x[i]).as_long() for i in range(3)]
    best_val = m.eval(objective).as_long()
    print(best_x)    # [1, 0, 0]
    print(best_val)  # 5
\end{lstlisting}
\end{example}

У овом примеру решење $\hat{\mathbf{x}}$ јесте легално, али није
оптимално: Z3 налази бољи дозвољени битски низ и тиме пружа егзактан
класични референтни резултат за поређење.

\begin{remark}
У пракси је ово један од најважнијих хибридних образаца:
QUBO се користи као формат за оптимизацију, а Z3 као формални слој
који доказује да су penalty чланови коректни, да су коефицијенти
довољни и да добијени битски низ заиста представља важеће решење
полазног проблема.
\end{remark}

% ============================================================
%  Поглавље 8: Закључак
% ============================================================
\section{Закључак}
\label{sec:zakljucak}

% --- Резиме доприноса рада ---
% Сумирај главне тачке из сваког поглавља у 2--3 реченице.

У овом раду разматрали смо НП-тешке проблеме кроз призму проблема
MaxCut и показали два комплементарна приступа: прецизно решавање
путем SMT решавача Z3, и апроксимативно решавање путем квантног
алгоритма QAOA.

Кључни закључци:

\begin{itemize}[leftmargin=1.5em]
  \item НП-тешки проблеми представљају суштинску препреку класичном
        рачунарству, а MaxCut је репрезентативни пример са богатом
        структуром применљивом на квантне технике.
  \item Z3 и SMT решавачи пружају прецизна решења за мале инстанце
        и служе као референца за верификацију квантних резултата.
  \item QAOA је природни квантни алгоритам за QUBO/Изингове проблеме;
        његова перформанса расте са дубином $p$, али га NISQ ограничења
        ограничавају на умерене дубине.
  \item Квантно рачунарство и аутоматско резоновање нису раздвојене
        области: хибридне петље, Гроверово убрзање и формална
        верификација квантних програма чине их природним саговорницима.
\end{itemize}

\subsection*{Отворена питања}

% --- Правци будућег истраживања ---
% Набројати 3--5 конкретних отворених проблема и праваца истраживања.

\begin{itemize}[leftmargin=1.5em]
  \item \textbf{Дубљи QAOA:} Да ли QAOA са великим $p$ може да
        надмаши Гоеманс-Вилијамсон апроксимацију ($0.878$) на
        свим класама графова?
  \item \textbf{Грешке и корекција:} Кад ће квантни хардвер са
        корекцијом грешака омогућити поуздани QAOA на великим
        инстанцама ($n > 1000$)?
  \item \textbf{Квантна верификација:} Развој аутоматских алата
        (на бази SMT/ATP) за верификацију исправности квантних
        алгоритама и кола.
  \item \textbf{Квантне SMT теорије:} Интеграција квантних оракла
        директно у SMT оквире као нове теорије (нпр. ``теорија
        квантне оптимизације'').
  \item \textbf{Хибридни решавачи:} Систематска истраживања
        хибридних алгоритама где Z3 (или CDCL) води квантно
        истраживање простора решења.
\end{itemize}

% TODO: Завршити закључак личним освртом на значај теме и
% потенцијалним утицајем квантних рачунара на будућност
% формалне верификације и аутоматског резоновања.


% ============================================================
%  ЛИТЕРАТУРА
% ============================================================
\begin{thebibliography}{9}

\bibitem{karp1972}
R.~M. Karp.
\newblock Reducibility among combinatorial problems.
\newblock In \textit{Complexity of Computer Computations}, pages 85--103.
  Plenum Press, 1972.

\bibitem{goemans1995}
M.~X. Goemans and D.~P. Williamson.
\newblock Improved approximation algorithms for maximum cut and satisfiability
  problems using semidefinite programming.
\newblock \textit{Journal of the ACM}, 42(6):1115--1145, 1995.

\bibitem{farhi2014}
E.~Farhi, J.~Goldstone, and S.~Gutmann.
\newblock A quantum approximate optimization algorithm.
\newblock \textit{arXiv preprint arXiv:1411.4028}, 2014.

\bibitem{monasson1999}
R.~Monasson et~al.
\newblock Determining computational complexity from characteristic `phase
  transitions'.
\newblock \textit{Nature}, 400:133--137, 1999.

\bibitem{demoura2008}
L.~de~Moura and N.~Bj{\o}rner.
\newblock Z3: An efficient SMT solver.
\newblock In \textit{TACAS 2008}, LNCS 4963, pages 337--340. Springer, 2008.

\bibitem{preskill2018}
J.~Preskill.
\newblock Quantum computing in the NISQ era and beyond.
\newblock \textit{Quantum}, 2:79, 2018.

\bibitem{nielsen2010}
M.~A. Nielsen and I.~L. Chuang.
\newblock \textit{Quantum Computation and Quantum Information}.
\newblock Cambridge University Press, 10th anniversary edition, 2010.

\end{thebibliography}

\end{document}
